\documentclass[a4paper,11pt]{article}
\usepackage{amsmath,amsthm,amsfonts,amssymb,amscd,amstext,vmargin,graphics,graphicx,tabularx,multicol} 
\usepackage[francais]{babel}
\usepackage[utf8]{inputenc}  
\usepackage[T1]{fontenc} 
\usepackage{pstricks-add,tikz,tkz-tab,variations}
\usepackage[autolanguage,np]{numprint} 

\setmarginsrb{1.5cm}{0.5cm}{1cm}{0.5cm}{0cm}{0cm}{0cm}{0cm} %Gauche, haut, droite, haut
\newcounter{numexo}
\newcommand{\exo}[1]{\stepcounter{numexo}\noindent{\bf Exercice~\thenumexo} : \marginpar{\hfill /#1}}
\reversemarginpar


\newcounter{enumtabi}
\newcounter{enumtaba}
\newcommand{\q}{\stepcounter{enumtabi} \theenumtabi.  }
\newcommand{\qa}{\stepcounter{enumtaba} (\alph{enumtaba}) }
\newcommand{\initq}{\setcounter{enumtabi}{0}}
\newcommand{\initqa}{\setcounter{enumtaba}{0}}

\newcommand{\be}{\begin{enumerate}}
\newcommand{\ee}{\end{enumerate}}
\newcommand{\bi}{\begin{itemize}}
\newcommand{\ei}{\end{itemize}}
\newcommand{\bp}{\begin{pspicture*}}
\newcommand{\ep}{\end{pspicture*}}
\newcommand{\bt}{\begin{tabular}}
\newcommand{\et}{\end{tabular}}
\renewcommand{\tabularxcolumn}[1]{>{\centering}m{#1}} %(colonne m{} centrée, au lieu de p par défault) 
\newcommand{\tnl}{\tabularnewline}

\newcommand{\bmul}[1]{\begin{multicols}{#1}}
\newcommand{\emul}{\end{multicols}}

\newcommand{\trait}{\noindent \rule{\linewidth}{0.2mm}}
\newcommand{\hs}[1]{\hspace{#1}}
\newcommand{\vs}[1]{\vspace{#1}}

\newcommand{\N}{\mathbb{N}}
\newcommand{\Z}{\mathbb{Z}}
\newcommand{\R}{\mathbb{R}}
\newcommand{\C}{\mathbb{C}}
\newcommand{\Dcal}{\mathcal{D}}
\newcommand{\Ccal}{\mathcal{C}}
\newcommand{\mc}{\mathcal}

\newcommand{\vect}[1]{\overrightarrow{#1}}
\newcommand{\ds}{\displaystyle}
\newcommand{\eq}{\quad \Leftrightarrow \quad}
\newcommand{\vecti}{\vec{\imath}}
\newcommand{\vectj}{\vec{\jmath}}
\newcommand{\Oij}{(O;\vec{\imath}, \vec{\jmath})}
\newcommand{\OIJ}{(O;I,J)}


\newcommand{\reponse}[1][1]{%
\multido{}{#1}{\makebox[\linewidth]{\rule[0pt]{0pt}{20pt}\dotfill}
}}

\newcommand{\titre}[5] 
% #1: titre #2: haut gauche #3: bas gauche #4: haut droite #5: bas droite
{
\noindent #2 \hfill #4 \\
#3 \hfill #5

\vspace{-1.6cm}

\begin{center}\rule{6cm}{0.5mm}\end{center}
\vspace{0.2cm}
\begin{center}{\large{\textbf{#1}}}\end{center}
\begin{center}\rule{6cm}{0.5mm}\end{center}
}



\begin{document}
\pagestyle{empty}
\titre{Interrogation: Fractions }{Nom :}{Prénom :}{Classe}{Date}


\exo{2} Compléter les phrases suivanates:

\bmul{2}
– Le double de 4 est . . . . . . . . . . . . . . .\\

– La moitie de 620 est . . . . . . . . . . . . .\\

– 60 est le . . . . . . . . . . . . . . . . . . .de 30.\\

– Le triple de 7 est . . . . . . . . . . . . . . . .\\



\columnbreak

– 2000 est le . . . . . . . . . . . . . . de 8000.\\

– Le tiers de 6 est . . . . . . . . . . . . . . . . .\\

– 21 est le . . . . . . . . . . . . . . . . . . . . de 7.\\

– 16 est la . . . . . . . . . . . . . . . . . . .de 32.
\emul




\exo{2,5} Pour chaque question, entourer \textbf{la} bonne réponse.\\

\includegraphics[scale=1]{fraction1.eps} \\


\exo{2} Compléter les égalités suivantes :

\bmul{4}

$\dfrac{1}{6}=\dfrac{.....}{30}$

\columnbreak

$\dfrac{35}{42}=\dfrac{.....}{6}$

\columnbreak

$\dfrac{3}{72}=\dfrac{1}{....}$

\columnbreak

$\dfrac{5}{7}=\dfrac{40}{.....}$\\

\emul



\exo{1,5} Donner les fractions irréductibles des fractions suivantes (simplifier au maximum ces fractions):

\bmul{4}
$E =\dfrac{42}{63}$\\

$E = ..............$\\

$E = ..............$

\columnbreak

$E = ..............$\\

$E = ..............$\\

$E = ..............$

\columnbreak

$K =\dfrac{60}{144}$\\

$K = ..............$\\

$K = ..............$

\columnbreak

$K = ..............$\\

$K = ..............$\\

$K = ..............$

\emul


\exo{2}

Dans une classe de 6ème de 24 élèves, le quart des élèves vient à pied et 512  des élèves viennent à vélo. Les autres se font amener au collège en voiture par leurs parents.\\

\q	Calculer le nombre d'élèves venant au collège à pied, à vélo et en voiture.\\
\reponse[4]\\

\q	Dans cette classe, quelle fraction des élèves vient au collège en voiture ?\\
\reponse[2]\\









\end{document}
